\documentclass{article}
\usepackage{amsmath}
\title{A qualified observation is useful only through the consuming cut}
\date{}
\begin{document}
\maketitle

\paragraph{Boundary.}
This is a composed, held-in civil logistics reference. The exact
native fragments and the shared admission core are pinned in
\verb|SOURCE.tsv|. The evidence graph and virtual transport seams
come from the earlier shared-stage replay. The old due time, clearance
predicate, shared transmitter map and consuming order are authored
interfaces. This is not a performance measurement, an independently
picked topology, a physical bypass traversal, or the stack's
confirmatory evaluation.

\paragraph{Non-waivable old work.}
At outer tick four, a native source latches whether the bay is clear
and sends the labeled sample to the six logical caches. A previous
failed send can occupy a stage transmitter so that several cache
refreshes fail together. Each cache can hold a genuinely acquired
prior reading with either value. Qualified reads additionally report
a final callback on this latch. Each serialized read takes one tick
by the declared service model. The new dispatch is due by the end
of tick six. The old carrier's authorized movement is due at tick
$r\in\{6,7,8\}$; it physically enters an initially clear bay or
leaves an occupied one then. There is no permission to stall it for
a late check. Its radio transfer at $r$ uses four bytes from Bypass
if initially clear and from Direct if initially occupied. The new
transfer must use one of the very same transmitters, but a distinct
receiver, and Direct additionally requires physical clearance at
its consuming cut. The equal-cut order is old movement, old
transfer, new transfer. Both keep their ordinary native callbacks,
and independent motion and traffic are not blocked.

\paragraph{Effect law and impossibility.}
A shared sender cannot pay both transfers: every tested cut where
new service is still possible is within the inclusive four-tick
window of old transfer, in either temporal order. Thus with initial
clearance the only feasible choice is Direct before $r$; with
initial occupancy the only feasible choice is Bypass, including
after $r$ while the due time still permits. These laws describe
guaranteed joint service, not merely a successful new callback.
Two independent executable checks, the common admission engine on
the joined events and a direct formula, agree for every bit, time,
route and old-work cut. In native controls, a wrong sender can make
the new callback succeed and the old one fail; at $r$, perfect
acquisition-time clearance can coexist with an unlawful actual
Direct callback even when the old job also succeeds.

For one admitted live debit, every label can still be stale and
the old cached bit can match either initial state. Two bare reads,
even chosen adaptively, cannot always separate the two states:
the exact earlier bare-read minimax bound is five, checked anew
here against the actionable decision. With authenticated
latch qualification a fixed pair $a,c$ always includes a refreshed
label, so the two reads suffice at $r\geq7$. If $r=6$, there is
an initially clear history for every first read in which that
cache is stale. Its first response also fits an occupied history,
so no single route is safe then. Even a perfect second read,
when it becomes available, cannot undo the required old movement.
This obstruction is separate from the two-debit case, where two
reads no longer suffice to guarantee a qualified response even
under the longer source lease.

\paragraph{Executable comparison.}
The complete finite search uses only each acquired local view.
At a decision node it tests all still-possible worlds, trying the
two explicit routes and all admissible unread labels. Withholding
preserves the old duties when no route can be justified. The
tree maximizes the count of distinct experiment rows served, which
is not a statistical estimate or a claim of optimizing frequency
of repeated execution fixtures. The full-service verdict is exact
because a branch sum reaches all rows iff every child is served.
The same-information direct search returns the same tree and
verdict. We run 16,896 executions across the 12 configurations
and preserve the old action and old transfer in every tested
execution of these controllers. At $r=7$ with one debit, the
qualified controller pays the new work in all 768 cases; the
chosen bare controller pays in 512. At $r=6$ qualification pays
in 640, but no policy in the qualified interface guarantees all.
These are finite sweep counts, not probabilities. The only strict
comparison uses a genuinely different input interface: the fresh
bit is callback-derived in the test, not provided by the stock
remote interface, and its acquisition cost and authenticity in an
external system remain open. The transfer guarantee is meaningful
only if the local time bound protects the same physical cut that
issued the clearance, and the old action ordering and sender debt
premises still hold.
\end{document}
